\documentclass[10pt,a4paper]{amsart}
\usepackage[margin=19mm]{geometry}
\usepackage{amsmath,amssymb,mathtools}
\usepackage[scr=rsfs,cal=euler]{mathalfa}
\usepackage{xfrac}
\usepackage[unicode,hidelinks,bookmarks=false]{hyperref}
\newcommand{\ie}{i.e.\ }
\newcommand{\cf}{cf.\ }
\newcommand{\resp}{respectively\ }
\newcommand{\Subsection}{\subsection}
\providecommand{\nobreakdash}{\nobreak\hbox{-}}
\setlength{\emergencystretch}{2em}
\multlinegap0pt
\usepackage{bm}
\usepackage{bbm}
\usepackage{dsfont}
\usepackage{xspace}
\usepackage{cancel}
\usepackage{mathtools}
\usepackage[shortlabels]{enumitem}
\usepackage{amsthm}
\makeatletter
\@ifundefined{corollary}{\newtheorem{corollary}{Corollary}}{}
\@ifundefined{lemma}{\newtheorem{lemma}{Lemma}}{}
\makeatother
\newcommand{\Lines}{\mathcal{L}}
\newcommand{\N}{\mathbb{N}}
\newcommand{\Points}{\mathcal{P}}
\title[On certain blocking sets and the minimum weight of the code ]{On certain blocking sets and the minimum weight of the code of generalised polygons}
\author{Sebastian Petit, Geertrui Van de Voorde}
\date{Source: arXiv:2511.07697v1 (CC BY 4.0)}
\begin{document}
\maketitle

\noindent\emph{Source and licence.} On certain blocking sets and the minimum weight of the code of generalised polygons. Version arXiv:2511.07697v1 (CC BY 4.0), distributed under the Creative Commons Attribution 4.0 International licence, as recorded in the arXiv licence metadata for this exact version and shown on the abstract page https://arxiv.org/abs/2511.07697v1. Full rights notice: licenses/arxiv-2511.07697-CC-BY-4.0.txt.\par\smallskip

\noindent\textit{Editorial scope.} Counted unit: the Lemma whose source label is \texttt{starLemma} in the article (source0.tex lines 402--408, source label \texttt{starLemma}), delivered together with the article's own complete original proof of it (source0.tex lines 409--442, one \texttt{proof} environment of 34 lines). No mathematics is written, repaired or re-derived: statement and proof are the article's own text line for line. Required context retained uncounted: intersections2 (lines 354--371). Every definition, display and label of those lines is retained unchanged. Omitted from this package and not redistributed: 1--353, 372--401, 443--860. Nothing inside a retained excerpt is omitted or altered: every retained body is delivered exactly as the article's lines are written, so no omission receipt of any kind accompanies this package. Citations: the retained excerpts carry no citation command at all, so no bibliography entry of the article is transcribed and no reference is redistributed. Reference adaptation: none was needed and none was made; every reference inside the delivered unit resolves inside it, so no unresolved reference or citation is produced. Printed slips kept literal: no printed word, symbol, hypothesis or number of the counted statement or of its proof was altered. Layout and class: the article's own class is \texttt{article} with packages including geometry, graphicx, xcolor, amsmath, amssymb, amsthm, enumitem, hyperref, comment; the wrapper is the lane's pinned \texttt{amsart} editorial wrapper at 10pt on A4, which restates the theorem-like environments the retained text needs and adds the article's own macro definitions that the retained text needs (\texttt{\textbackslash Lines}, \texttt{\textbackslash N}, \texttt{\textbackslash Points}), with extra wrapper packages (bm, bbm, dsfont, xspace, cancel, mathtools, enumitem). The delivered numbering is the unit's own. Assets: no retained span contains a figure environment, a graphics-inclusion command, a PDF-page inclusion, a table or a TikZ picture; no image file is copied into this package, the package declares no asset dependency and no image file is redistributed with it. Licence: version arXiv:2511.07697v1 of this article is distributed under the Creative Commons Attribution 4.0 International licence, as recorded in the arXiv licence metadata for this exact version and shown on the abstract page https://arxiv.org/abs/2511.07697v1. A full rights notice is delivered at licenses/arxiv-2511.07697-CC-BY-4.0.txt. Characters: every retained excerpt is pure ASCII, so the delivered engine, XeLaTeX with its default Latin Modern text font, reports no missing character.
\begin{corollary}\label{intersections2}
\begin{enumerate}[(i)]
\item 	Let \(x\) be an element of $\Gamma$. Let \(y_1\neq y_2 \in \Gamma_1(x)\) and let \(d\in\{1,\dots,2m-2\}\).
	Then,
	\[\Gamma_{\le d - 1}(x) = \Gamma_{\le d}(y_1)\cap\Gamma_{\le d}(y_2).\]

\item  Let \(L\) be a line of \(\Gamma\) and let \(r\in \{1,2,\dots,m-1\}\).
	\begin{enumerate}[(a)]
		\item If 
		\(M_1\) and \(M_2\) are two disjoint lines intersecting \(L\),
		then
		\[\Points_{\le 2r-3}(L) = \Points_{\le 2r-1}(M_1)\cap\Points_{\le 2r-1}(M_2).\]
		\item If
		\(v_1\) and \(v_2\) are two distinct points on \(L\), then
		\[\Points_{\le 2r-1}(L) = \Points_{\le 2r}(v_1)\cap\Points_{\le 2r}(v_2).\]
	\end{enumerate}
\end{enumerate}
\end{corollary}

\begin{lemma}\label{starLemma}
	Let \(C\) be a set of points of $\Gamma$ with \(|C|<s+1\). Let \(d\in \{1,2,\dots,m-1\}\).
	\begin{enumerate}[(i)]
		\item \(\forall L \in \Lines: \exists M \in \Lines_{\le 2d-2}(L): C\cap \Points_{\le2d-1}(M) = C\cap \Points_{1}(L)\).
		\item \(\forall L \in \Lines: \exists v \in \Points_{\le 2d-1}(L): C\cap \Points_{\le2d}(v) = C\cap \Points_{1}(L)\).
	\end{enumerate}
\end{lemma}

\begin{proof}
	\(\)
	\begin{enumerate}[(i)]
		\item We prove this by induction on \(d\). The base case \(d=1\) follows by taking $M=L$.
		Assume that the statement holds for some \(k\in \N\) with \(1\le k < m-1\) and consider the case \(d=k+1\). 
		Let \(L\in\Lines\) be an arbitrary line of \(\Gamma\). Then there exists an \(M\in \Lines_{\le 2k-2}(L)\) such that \(C\cap \Points_{\le2k-1}(M) = C\cap \Points_{1}(L)\). 
		
		We label the \(s+1\) points on \(M\) as \(\{v_1,v_2,\dots,v_{s+1}\}\). For each of these points \(v_i\) we can consider a line \(N_i\) through \(v_i\) and different from \(M\). Note that the lines $N_1,\ldots,N_{s+1}$ are mutually disjoint.
		
		By Corollary \ref{intersections2}(2a), we now observe that for any two of these lines \(N_i\) and \(N_j\) with \(i\neq j\), we have \[\Points_{\le 2(k+1)-1}(N_i)\cap \Points_{\le 2(k+1)-1}(N_j)= \Points_{\le 2(k+1)-3}(M).\]
		
		Since  \(|C|<s+1\), at least one of the $s+1$ disjoint sets \(\Points_{\le 2(k+1)-1}(N_i)\setminus \Points_{\le 2(k+1)-3}(M) \), $i \in \{1,2,\ldots,s+1\}$, must have an empty intersection with \(C\). Without loss of generality, we may assume 
		\[(\Points_{\le 2(k+1)-1}(N_1)\setminus \Points_{\le 2(k+1)-3}(M)) \cap C = \emptyset.\]
		This means that 
		\[C\cap \Points_{\le2(k+1)-1}(N_1) = C\cap \Points_{\le2(k+1)-3}(M) = C\cap \Points_{1}(L),\]
		 where we have invoked the induction hypothesis.
         
		Since \(M\in \Lines_{2k-2}(L)\) and \(N_1\in\Lines_{2}(M)\),  we see that \(N_1\in   \Lines_{2k-2+2}(L)=\Lines_{2k}(L)=\Lines_{2(k+1)-2}(L)\). The statement follows.
		
		\item Let \(L\in\Lines\) be an arbitrary line of \(\Gamma\). 
		From part (i), we know that there exists a line \(M\in \Lines_{\le 2d-2}(L)\) such that \(C\cap \Points_{\le2d-1}(M) = C\cap \Points_{1}(L)\).
		
		We label the \(s+1\) points on \(M\) as \(\{v_1,v_2,\dots,v_{s+1}\}\). 

		By Corollary \ref{intersections2}(2b), we have that for \(i\neq j\),  \[\Points_{\le 2d}(v_i)\cap \Points_{\le 2d}(v_j)= \Points_{\le 2d-1}(M).\]
		
		Since \(|C|<s+1\) points, at least one of the disjoint  sets \(\Points_{\le 2d}(v_i)\setminus \Points_{\le 2d-1}(M) \) must have an empty intersection with \(C\). Without loss of generality, we may assume 
		\[(\Points_{\le 2d}(v_1)\setminus \Points_{\le 2d-1}(M)) \cap C = \emptyset.\]
		This means that 
		\[C\cap \Points_{\le2d}(v_1) = C\cap \Points_{\le2d-1}(M) = C\cap \Points_{1}(L)\]
		
		Since $v_1\in \Points_1(M)$ and $M\subseteq \Lines_{\leq 2d-2}(L)$, we see that \(v_1\in \Points_{1}(M) \subseteq \Points_{\leq 2d-1}(L)\) and the statement follows.\qedhere
	\end{enumerate}
\end{proof}

\end{document}
